\section{Measurement}\label{sec:apparatus}
The detector is a tracker of \InputsDetectorPlanes{} scintillator planes with \InputsBarsPerPlane{} bars each. Its active width is $\InputsActiveWidthCm\,\mathrm{cm}$ and the gap between the outer layers is $\InputsLayerGapCm\,\mathrm{cm}$. We took data at two floor positions under the ceiling (pos0 and pos1) and on the roof with open sky above, using the same detector for all three runs.

After selection we kept \DataFirstTracks{} tracks in $\DataFirstHours\,\mathrm{h}$ at pos0, \DataSecondTracks{} tracks in $\DataSecondHours\,\mathrm{h}$ at pos1 and \DataSkyTracks{} tracks in $\DataSkyHours\,\mathrm{h}$ under open sky. Live time was taken from the distribution of time intervals between tracks. The acquisition files do not state the unit of these intervals; we inferred it from the trigger rate.

The origin is at pos0, $z$ points up and the main beams run along $y$. The position and orientation of pos1 relative to pos0 were fitted from the muon data (Sec.~\ref{sec:method}); only a rough floor distance was known. On site, the beam bottom is at $\InputsMeasuredBeamBottom\,\mathrm{m}$ and the vertical beam thickness is $\InputsMeasuredBeamDepth\,\mathrm{m}$, from a rough tape measurement without stated uncertainty. The detector reference height above the floor was not calibrated, so direct height comparison assumes a common origin.

\begin{figure}[t]
\centering
\includegraphics[width=\columnwidth]{setup.pdf}
\caption{Measurement layout and beam array. Beams beyond the matched set are drawn by continuing the pitch; box sizes are illustrative.}\label{fig:setup}
\end{figure}
